% !TeX root = main_2.tex

\documentclass[11pt,fleqn]{article}
\usepackage{uebungen}
\renewcommand{\logo}{\includegraphics[height=10ex]{kit_logo_de_1c_schwarz}}

\usepackage{amsfonts}
\usepackage[dvipsnames]{xcolor}
\usepackage[hidelinks,colorlinks=true,linktocpage=true,linkcolor=NavyBlue,urlcolor=NavyBlue,citecolor=NavyBlue]{hyperref}
%\usepackage[hang]{subfigure}
\usepackage{lmodern}
%\usepackage[scaled=.85]{beramono}
\usepackage{listings}

\usepackage{cleveref}

%\usepackage{caption}
\usepackage{subcaption}
\usepackage{graphicx}

\usepackage{amsmath}
\usepackage{amsthm}
\usepackage{amssymb}
\newtheorem{definition}{Definition}
\newtheorem{lemma}{Lemma}

\usepackage{todonotes}

\let\vec\mathbf

\definecolor{shellBack}{rgb}{0.95,0.95,0.9}
\lstset{basicstyle=\ttfamily, backgroundcolor=\color{shellBack}}

\def\published{true}

\begin{document}

\kopfBlatt{Neural Network Verification (Exercises 2)}

\section*{Incomplete Verification using Zonotopes}

On the last exercise sheet, we used mixed integer linear programming (via the \texttt{MIPVerify} solver) to verify 
properties of image classifiers and airborne collision avoidance systems (ACAS).

The \texttt{DNNV} framework also provides a zonotope-based verifier called \texttt{DeepZono}, which can be called via the following command:
\begin{verbatim}
    dnnv -v --seed 42 -N N <network.onnx> --eran <property.py>
\end{verbatim}
(make sure to include \texttt{--seed 42} to ensure reproducibility).
If the property you want to prove has a defines a parameter via e.g. 
\begin{verbatim}
    epsilon = Parameter("epsilon", default=1/255)
\end{verbatim}
you can pass a value for this parameter via the command line as follows:
\begin{verbatim}
    dnnv -v -N N <network.onnx>  --eran --prop.epsilon=0.02 <property.py
\end{verbatim}

In this exercise, we will see, how incomplete zonotope-based verifiers compare to the MIP-based approach.

\paragraph{Task 1:}
Use the \texttt{DeepZono} verifier to verify the properties from the previous exercise sheet.
How does \texttt{DeepZono} compare to \texttt{MIPVerify} with respect to runtime?
Are there instances that can be solved by \texttt{MIPVerify} but not by \texttt{DeepZono}?


\section*{Confidence-based Equivalence Verification}

When neural networks are used in resource-constrained or performance-critical environements,
they are often pruned to reduce their size and the number of operations required to evaluate them.
However, we still want to ensure that the pruned network behaves similarly to the original one.

In this exercise, we are going to look at a network trained to classify particla jets in CERN's 
Large Hadron Collider (LHC) experiments.
In this setting, large amounts of data have to be processed in real-time, making pruning highly desirable.

Specifically, we will look at the \texttt{lhc\_2\_20-1\_no\_softmax.onnx} network, and a compressed version of it.
In the \texttt{lhc\_2\_20-1\_0.1\_no\_softmax.onnx} network, 10\% of the weights have been set to zero.

As neither \texttt{MIPVerify} nor \texttt{DeepZono} support the softmax activation function, we removed it from the networks.
However, remember that we can overapproximate it for a confidence $\delta \geq \frac{1}{2}$ by the following linear constraints: 
\begin{align*}
    \left\{
        x \in \mathbb{R}^m~\middle|~
        \mathrm{softmax}_i\left(x\right) \geq \delta
    \right\}
\subseteq
\left\{
x \in \mathbb{R}^m~\middle|~
x_i \geq x_j + \ln\left(\frac{\delta}{1-\delta}\right)
\right\}
\end{align*}

\paragraph{Task 2.1:}
The \texttt{conf\_equiv\_0.py} property file already contains code to constrain the inputs of the neural 
networks to an $\varepsilon$-box around the origin.
Complete \texttt{conf\_equiv\_0.py} such that it checks the following property:
\\
\\
\emph{If the original network classifies its inputs as a gluon (output class $0$), with a confidence of at leat $\delta$, 
then the compressed network also classifies the input as a gluon.}
\\
\\
\textit{Hint 1:} 
You can access the $j$-th output of a neural network \texttt{N} given an input $x$ via \texttt{N(x)[0,j]} (the first index is the batch dimension).

\paragraph{Task 2.2:}
Try to verify the property using \texttt{DeepZono}.
How large can you set $\varepsilon$ for a specific $\delta$ such that the property can still be verified?
How large can you set $\delta$ for a specific $\varepsilon$ such that the property can still be verified?


\section*{Scalability}

Verification of $\mathrm{ReLU}$ neural networks is a decidable, but NP complete problem.
A natural question is, how well the verification algorithms scale in practice and which
parameters influence verification runtime.

To this end, we will look at a range of targeted robustness properties
\begin{equation*}
    \forall \vec{x} \in [\vec{x}_0 - \varepsilon, \vec{x}_0 + \varepsilon]: f_{\mathcal{N}}(\vec{x})_0 \geq f_{\mathcal{N}}(\vec{x})_1
\end{equation*}
parameterized by different networks $\mathcal{N}$ and values for $\varepsilon$.

\paragraph{Task 3:}
Run \texttt{MIPVerify} and \texttt{DeepZono} with different combinations of 
\begin{itemize}
    \item networks: \texttt{mnist\_relu\_3\_50.onnx} or \texttt{mnist\_relu\_9\_200.onnx}
    \item perturbation radius $\varepsilon$
    \item bounds tightening algorithm for \texttt{MIPVerify}: Interval arithmetic or linear programming
\end{itemize}
You can run the property using the following commands for \texttt{MIPVerify} and \texttt{DeepZono}:
\begin{verbatim}
    dnnv -v --seed 42 -N N <network.onnx> --mipverify --mipverify.optimizer GLPK 
            --mipverify.tightening [interval_arithmetic|lp] 
            --prop.epsilon=<...> targeted_robustness.py

    dnnv -v --seed 42 -N N <network.onnx> --eran 
            --prop.epsilon=<...> targeted_robustness.py
\end{verbatim}
How do the parameters or the network affect the runtime of the verifiers?
\\
\\
\textit{Hint:}
\''Interesting\'' combinations of values are shown in Table \ref{tab:interesting_combinations}.

\begin{table*}
\centering
\begin{tabular}{|c|c|c|c|c|}
\hline
    Network & $\varepsilon$ & MIPVerify Tightening \\
\hline
    mnist\_relu\_3\_50.onnx & 0.046  & IA  \\  % both unsat, MIP 30sec, DeepZono 1.1sec
    mnist\_relu\_3\_50.onnx & 0.047 & IA / LP  \\  % both unsat, MIP (IA) 150sec, MIP (LP) 28sec, DeepZono 1.2sec
    mnist\_relu\_3\_50.onnx & 0.06 & LP \\  % MIP unsat 64sec, DeepZono unknown 1.2sec
    mnist\_relu\_9\_200.onnx & 0.0039 & IA \\  % both unsat, MIP 28sec, DeepZono 3sec
    mnist\_relu\_9\_200.onnx & 0.0078 & IA \\  % both unsat, MIP 29sec, DeepZono 3sec
    mnist\_relu\_9\_200.onnx & 0.016 & LP \\  % MIP unsat 133sec, DeepZono unknown 4sec
\hline
\end{tabular}
\caption{Interesting combinations of networks, $\varepsilon$, and verifiers.
         The \texttt{DeepZono} verifier can be run for each of the above entries.
         \texttt{MIPVerify} can be run with the \texttt{--mipverify.tightening} flag set to either \texttt{interval\_arithmetic} or \texttt{lp} as 
         indicated in the last column.}
\label{tab:interesting_combinations}
\end{table*}






\end{document}